%%%%%%%%%%%%%%%%%%%%%%%%%%%%%%%%%%%%%%%%%%%%%%%%%%%%%%%%%%%%%%
%%%%%%%%%%%%% MA-XXXX Introducción al Cálculo en una variable %%%%%%%%%%%%%%%%
%%%%%%%%%%%%%%%%%%%%%%%%%%%%%%%%%%%%%%%%%%%%%%%%%%%%%%%%%%%%%%
\newpage
\paragraph*{MA-02XX Introducción al Cálculo en una variable}
\addcontentsline{toc}{subsection}{MA-02XX Introducción al Cálculo en una variable}

\begin{refsection}
\begin{table}[H]
\centering{
\begin{tabular}{|ll|}
\hline
%\multicolumn{2}{|c|}{\textbf{Nivel de virtualidad del curso:} Virtual}\\ \hline
\textbf{Año:} I  & \textbf{Requisitos:} Ninguno \\ 
\textbf{Ciclo:} I  & \textbf{Correquisitos:} Ninguno \\ 
\textbf{Tipo de curso:} Teórico & \textbf{Horas presenciales por semana:}   \\ 
\textbf{Créditos:} 3 & \textbf{Horas de trabajo independiente por semana:}  \\ \hline
\end{tabular}
}
\end{table}



\subsubsection*{Descripción del curso}

El Cálculo Diferencial e Integral constituye uno de los más grandes triunfos intelectuales de la historia de la humanidad. Este ha posibilitado una gran parte de los avances tecnológicos y científicos que se han tenido en los últimos siglos. Además, su descubrimiento a finales del siglo 17 por Newton y Leibniz  abrió senderos en el campo de la matemática que todavía hoy se están explorando.
Esta área del conocimiento, y los conceptos que introduce, van a aparecer repetidamente durante el resto de la carrera.

En este primer curso de Cálculo, se presentan los conceptos fundamentales desde un punto de vista concreto, enfatizando el entendimiento por medio de la práctica. En particular, se empezará el curso estudiando el eje unificador del cálculo, que es el concepto de límite de una función. A partir de esto se pasará a estudiar nociones básicas como continuidad, el concepto de derivada y de integral definida de una función, todas estas definidas en términos de límites. Uno de los objetivos fundamentales del curso será desarrollar una variedad de técnicas para el cálculo de límites, derivadas e integrales.

Este es un curso introductorio al cálculo diferencial e integral en una variable con un énfasis en el desarrollo de habilidades operacionales y del pensamiento intuitivo y geométrico con base en la exploración de una serie de conceptos y procedimientos propios del cálculo. Para lograr el desarrollo de estas habilidades e intuiciones, se priorizará el trabajo con ejemplos concretos.

Este curso se complementa con MA-1XXX Matemática Exploratoria en el desarrollo de la intuición sobre objetos básicos de la matemática que serán primordiales en la construcción y estudio de estructuras con un mayor nivel de complejidad o abstracción en cursos posteriores de la carrera.




\subsubsection*{Objetivos}

\textbf{General:} Aplicar las nociones, intuiciones y conceptos básicos del cálculo diferencial e integral en una variable para la resolución de problemas. \\

\textbf{Específicos:}

\begin{enumerate}
\item Definir de manera intuitiva los conceptos de funciones, límites, derivadas e integrales.
\item Calcular límites, derivadas e integrales de funciones de una variable real.
\item Interpretar gráficamente los conceptos de funciones, límites, derivadas e integrales.
\item Relacionar los conceptos de derivada e integral como procesos inversos.
\item Aplicar los conceptos de funciones, límites, derivadas e integrales para la resolución de problemas concretos.
\item Demostrar existencia de límites de funciones concretas mediante la definición con $\varepsilon$ y $\delta$.
\item Indagar en la literatura sobre temas que permitan complementar su formación.
\item Comunicar ideas matemáticas de manera clara y efectiva en forma oral y escrita.
\end{enumerate}

\subsubsection*{Contenidos}

\begin{enumerate}
\item \textbf{Repaso de álgebra elemental (1 semana):} 
\begin{enumerate}
    \item Potencias.
    \item Fórmulas notables.
    \item Factorización.
    \item Racionalización.
    \item Valor absoluto y distancia.
    \item Desigualdades.
\end{enumerate}

\item \textbf{Introducción a las funciones (1 semana):} 
\begin{enumerate}
    \item Repaso básico de gráficas de funciones elementales.
    \item Propiedades básicas de funciones (monotonía, paridad, inyectividad, sobreyectividad, periodicidad)
    \item Funciones inversas.
    \item Transformaciones básicas (reflexiones, traslaciones).
\end{enumerate}

\item \textbf{Límites y continuidad (3 semanas):} 
\begin{enumerate}
    \item Concepto intuitivo de límite.
    \item Propiedades básicas de límites (álgebra de límites, Teoremas del Encaje).
    \item Cálculo de límites de funciones elementales (funciones polinomiales, racionales, con radicales).
    \item Continuidad (teoremas básicos, composición de funciones continuas, cambios de variable, Teorema de los Valores Intermedios).
    \item Límites infinitos y al infinito. Asíntotas horizontales y verticales.
    \item Límites trigonométricos básicos.
\end{enumerate}

\item \textbf{Derivación (5 semanas):} 
\begin{enumerate}
    \item Concepto geométrico y físico de la derivada.
    \item Cálculo de derivadas con la definición.
    \item Reglas básicas de derivación (suma, producto, cociente, regla de cadena, derivada de la inversa, derivadas de orden superior).
    \item Derivación implícita y logarítmica.
    \item Teorema del Valor Medio y sus aplicaciones.
    \item Aproximaciones de Taylor. 
    \item Regla de L'Hopital. Aplicación a límites con exponenciales y logaritmos. 
    \item Graficación de funciones con técnicas de cálculo.
    \item Optimización.
\end{enumerate}

\item \textbf{Integración de Riemann (4 semanas):} 
\begin{enumerate}
    \item La integral definida como el área bajo una curva.
    \item Idea intuitiva de la integral definida como el límite de sumas de Riemann. Cálculo de ejemplos concretos con polinomios de grado a lo sumo 3.
    \item Antiderivadas y el Teorema Fundamental del Cálculo.
    \item Reglas de integración básicas (sustitución, integración por partes, fracciones parciales, integración de funciones trigonométricas, sustitución trigonométrica).
\end{enumerate}

\item \textbf{Introducción al concepto formal de límite de una función (1 semana):} 
\begin{enumerate}
    \item Definición formal de límite con $\varepsilon$-$\delta$.
    \item Interpretación geométrica de la definición formal de límite.
    \item Demostración de límites concretos usando la definición formal. 
\end{enumerate}

\end{enumerate}

\subsubsection*{Metodología}

\noindent \textbf{Lineamientos metodológicos:} La persona docente a cargo de este curso debe respetar las siguientes pautas:
\begin{enumerate}
    \item La presentación de los contenidos debe priorizar la intuición y el desarrollo de habilidades operacionales por encima del formalismo. Por ejemplo se deben evitar las demostraciones que dependan del axioma del extremo superior o sus equivalentes, que se estudiarán en un curso posterior de principios de Análisis Matemático.
    \item Los argumentos con $\varepsilon$-$\delta$ se reservarán para el último tema, en el cual estos se aplicarán únicamente con ejemplos concretos.
    \item Las demostraciones que se deben realizar en este curso son las que contribuyen al desarrollo de la intuición y las habilidades operacionales en la persona estudiante.
    %\item Como existe codependencia entre la temática y objetivos de este curso y la del curso MA-1XXX Précalculo para Matemáticos, debe respetarse estrictamente el orden de temas establecido en este documento.
\end{enumerate}

\textbf{Sugerencias metodológicas:} Adicionalmente, se tienen las siguientes recomendaciones para la persona docente a cargo de este curso:
\begin{enumerate}
    \item Utilizar el recurso tecnológico para reforzar distintos conceptos estudiados en el curso por medio de la visualización, cálculo y experimentación. Se pone a disposición de la persona docente un repositorio de herramientas listas para ser implementadas en el curso, tanto en las clases como para el trabajo independiente de las personas estudiantes. Por ejemplo, al tratar algunos límites o derivadas de funciones trascendentes se pueden hacer gráficas o tablas de valores ya que en cursos posteriores se dará un tratamiento formal de esto.
    \item Aprovechar momentos adecuados del curso para motivar el estudio por medio de reseñas acerca del desarrollo histórico del cálculo y su importancia en aplicaciones dentro y fuera de la matemática
    \item En la evaluación, se sugiere utilizar estrategias que promuevan el trabajo constante de parte de las personas estudiantes a lo largo del ciclo lectivo. Por ejemplo, esto se puede hacer por medio de tareas o quizzes.
    \item Para fomentar la iniciación de las personas estudiantes en el desarrollo de sus habilidades investigativas, se sugiere la implementación de pequeños proyectos de investigación. Por ejemplo, pueden ser proyectos de investigación bibliográfica en temas históricos o también proyectos cortos de aplicación del cálculo como los que se encuentran en el libro de James Stewart \cite{Ste18}.
\end{enumerate}




\nocite{Abb15}
\nocite{Apo84}
\nocite{Apo85} 
\nocite{BS11}
\nocite{Apo76}
\nocite{SRW17}
\nocite{Ste18}
\nocite{Spi08}
\nocite{Ros13}

\printbibliography[heading=subbibliography]
\end{refsection}
